\chapter{Vitamin D Deficiency --- Rickets and Osteomalacia}
\index{Vitamin D Deficiency --- Rickets and Osteomalacia}
\label{sec:Vitamin D Deficiency}

\section{Overview}
Vitamin D deficiency is a vitamin deficiency disorder resulting from inadequate sunlight exposure, low dietary intake (or lack of fortified foods), malabsorption (celiac disease, cystic scarring, short bowel syndrome), kidney disease (impaired 1$\alpha$-hydroxylation), obesity (sequestration in adipose tissue), and medications (antiepileptics, rifampin, glucocorticoids). This metabolic disorder involves dysregulation of normal bodily processes, often requiring ongoing monitoring and management. It is increasingly common due to lifestyle and dietary factors. Metabolic disorders involve disruption of normal biochemical processes essential for health. These conditions may result from enzyme deficiencies, hormonal imbalances, or lifestyle factors. Many metabolic disorders are chronic and require ongoing management. Lifestyle modifications are often the cornerstone of treatment.
\section{Causes}
This condition happens because of deficiency of vitamin D, a fat-soluble secosteroid essential for intestinal calcium absorption, bone mineralization, and calcium homeostasis. Metabolic disorders arise from dysregulation of normal biochemical pathways. This may result from enzyme deficiencies, hormonal imbalances, or lifestyle factors that overwhelm normal homeostatic mechanisms. The underlying mechanisms involve complex interactions between genetic predisposition and environmental factors. Disruption of normal physiological processes leads to the characteristic features of the condition. Multiple contributing factors may act synergistically to trigger or worsen the condition.
\section{Risk Factors}
\begin{itemize}[noitemsep]
  \item Genetic predisposition and family history
  \item Advancing age
  \item Lifestyle factors (diet, exercise, smoking, alcohol)
  \item Environmental exposures
  \item Underlying medical conditions
  \item Medication use
\end{itemize}
\section{Signs and Symptoms}

\subsection{Common symptoms}
\begin{itemize}[noitemsep]
  \item Clinical features in children include rickets (defective mineralization of growing bone) with craniotabes, frontal bossing, delayed fontanelle closure, rachitic rosary, Harrison groove, bowing of the legs (genu varum), widened wrists and ankles, myopathy, and hypocalcemic tetany. In adults, osteomalacia presents with bone pain, proximal muscle weakness (waddling gait), and fragility fractures.
  \item Pain or discomfort in the affected area
  \end{itemize}

\subsection{Emergency warning signs}
\begin{itemize}[noitemsep]
  \item Severe or worsening symptoms of vitamin d deficiency --- rickets and osteomalacia require immediate medical evaluation
\end{itemize}
\section{Progression and Stages}
The condition may progress through different stages of severity over time. Early stages often involve mild or subtle symptoms that may be overlooked. Moderate stages involve more noticeable symptoms affecting daily function. Advanced stages may involve significant impairment and complications.
\section{Complications}
\begin{itemize}[noitemsep]
  \item Disease progression leading to worsening symptoms
  \item Secondary infections or injuries
  \item Side effects of treatment
  \item Reduced quality of life
  \item Emotional and psychological impact
\end{itemize}
\section{Diagnosis}
Doctors diagnose this using serum 25-hydroxyvitamin D less than 20 ng/mL (50 nmol/L), elevated alkaline phosphatase, low-normal calcium, and low phosphate

\begin{itemize}[noitemsep]
  \item Radiographs of rickets show cupping, fraying, and widening of the metaphysis
  \item In osteomalacia, Looser zones (pseudofractures) are seen.
  \item Comprehensive medical history and physical examination
  \item Laboratory tests to assess organ function
  \item Imaging studies to visualize affected structures
  \item Specialized testing depending on the affected system
  \item Consultation with relevant specialists
\end{itemize}
\section{Prevention}
\begin{itemize}[noitemsep]
  \item Maintain a healthy lifestyle with balanced diet and exercise
  \item Avoid known risk factors
  \item Attend regular medical check-ups for early detection
  \item Follow recommended screening guidelines
\end{itemize}
\section{Treatment Options}
\begin{itemize}[noitemsep]
  \item Treatment focuses on vitamin D 2000--5000 IU daily for 12 weeks or 50,000 IU weekly for 8 weeks, followed by maintenance 800--2000 IU daily
  \item Correction of calcium and phosphate is essential.
  \item Medications to manage symptoms and address underlying mechanisms
  \item Lifestyle modifications and self-care measures
  \item Physical therapy and rehabilitation
  \item Surgical intervention when indicated
  \item Regular monitoring and follow-up care
\end{itemize}
\section{Living with It}
\begin{itemize}[noitemsep]
  \item Follow your treatment plan as prescribed
  \item Maintain regular follow-up with your healthcare provider
  \item Make necessary lifestyle adjustments
  \item Seek support from family, friends, or support groups
  \item Stay informed about your condition
\end{itemize}
\section{Outlook}
Most people recover well with treatment.
